By Fact \claimref{synthesis.imported-gtg-dictionary}(ii) the rotation-gauge field $\Omega$ is the unique torsion-free connection of the gravitational field. Write $f^0 = \nabla\tau$, $f^i = \tau^{p_i}\nabla x^i$, so that $\hb^{T}(f^a) = \gamma^a$; on the constant physical frame the torsion equation of Definition \claimref{obj.synthesis.gtg-flat-spacetime-gauge-fields}(f) reads $\hb^{T}(\nabla\wedge f^a) = \sum_b \gamma^b\wedge(\omega(\gamma_b)\cdot\gamma^a)$. The left side: $\nabla\wedge f^0 = 0$, and $\nabla\wedge f^i = \nabla(\tau^{p_i})\wedge\nabla x^i = p_i\tau^{p_i-1}\nabla\tau\wedge\nabla x^i + \log\tau\,\tau^{p_i}\,\nabla p_i\wedge\nabla x^i$, whose image under the outermorphism $\hb^{T}$ is $(p_i/\tau)\,\gamma^0\wedge\gamma^i + \log\tau\,\hat\nabla p_i\wedge\gamma^i$. The right side for the displayed $\omega$: the boost terms give $(p_i/\tau)\,\gamma^0\wedge\gamma^i$ for $a = i$ and $\sum_k(p_k/\tau)\,\gamma^k\wedge\gamma_k = 0$ for $a = 0$, as in Proposition \claimref{dark-energy.frozen-point-datum-kasner-wall}; the tangential terms $-\log\tau\,\hat\nabla p_k\wedge\gamma_k$ are spacetime bivectors orthogonal to $\gamma^0$, so they give nothing for $a = 0$, and for $a = i$, with $(u\wedge\gamma_k)\cdot\gamma^i = \delta_{ki}\,u - (u\cdot\gamma^i)\,\gamma_k$, they give $-\log\tau\,\bigl(\gamma^i\wedge\hat\nabla p_i - \sum_k(\hat\nabla p_k\cdot\gamma^i)\,\gamma^k\wedge\gamma_k\bigr) = \log\tau\,\hat\nabla p_i\wedge\gamma^i$, since $\gamma^k\wedge\gamma_k = 0$. Both sides agree term by term, and uniqueness of the torsion-free connection makes this the Levi-Civita connection. (In components, a labelled reading: this is the first structure equation $de^i + \omega^{i}{}_{0}\wedge e^0 + \omega^{i}{}_{j}\wedge e^j = 0$ with $\omega^{i}{}_{j} = -\omega^{ij}$ in the convention of \eqref{eq:ec}, the tangential block being the Levi-Civita connection $\omega^{ij} = (\partial_i a_j/a_i)dx^j - (\partial_j a_i/a_j)dx^i$ of the diagonal three-metric $\sum_i a_i^2(dx^i)^2$ with $a_i = \tau^{p_i(x)}$.) Orders: $\hat\nabla p_i$ has components $\tau^{-p_j}\partial_j p_i$, so the tangential block is $O(\tau^{-p_{\max}}|\log\tau|)$ against $O(\tau^{-1})$ for the boost block, and $p_{\max} < 1$ on the open Kasner circle. In the field strength $R(a\wedge b) = a\cdot\nabla\Omega(b) - b\cdot\nabla\Omega(a) + \Omega(a)\times\Omega(b)$, read on the physical frame planes through $\mathcal{R}(B) = -R(\hb(B))$ (the same orders as the components of the curvature two-forms): the derivative of the boost block along $x^j$ contributes $(\partial_jp_i)(1 + p_i\log\tau)\tau^{-1-p_j}$ on the plane $\gamma_j\wedge\gamma_i$; the commutator of the rotation and boost blocks is $O(\tau^{-1-p_j}|\log\tau|)$; the derivatives of the rotation block are $O(\tau^{-1-p_j}|\log\tau|) + O(\tau^{-p_j-p_k}|\log\tau|)$; and the commutator of the rotation block with itself is $O(\tau^{-p_j-p_k}\log^2\tau)$. Against the Kasner terms of order $\tau^{-2}$, each is of relative order at most $\tau^{1-p_{\max}}|\log\tau|^2$, which tends to zero. Hence the Ricci map at order $\tau^{-2}$ is that of Proposition \claimref{dark-energy.frozen-point-datum-kasner-wall}(ii) with $p = p(x)$; the slice value of $Q$ at order $\varepsilon^{-2}$ is that of Proposition \claimref{dark-energy.wall-counterterm}(ii) with $p = p(x)$ by the same order counting (in that Proposition's components, the quadratic term of $Q$ no longer vanishes identically, since $d\omega^{ij}$ has tangential legs, but it is of lower order in $\tau^{-1}$); and the subleading orders carry the momentum constraints, which are not analysed here. The script confirms the displayed connection as an exact identity: with $\log\tau$ and the gradients $\partial_jp_i$ replaced by independent rational placeholders, the displayed $\Omega$ satisfies the torsion equation of the gravitational field and the torsion map has rank $24$, and the connection with the opposite tangential sign does not satisfy it.
